\subsection{OPEN SETS}

DEFINITION 1. A topological structure (or, more briefly, a topology) on a set X is a structure given by a set D of subsets of X, having the following properties (called axioms of topological structures):

\((\mathbf{O}_{\mathbf{I}})\) Every union of sets of \(\mathfrak{D}\) is a set of \(\mathfrak{D}\) .

\((\mathbf{O}_{\mathbf{II}})\) Every finite intersection of sets of \(\mathfrak{D}\) is a set of \(\mathfrak{D}\) .

The sets of \(\mathfrak{D}\) are called open sets of the topological structure defined by \(\mathfrak{D}\) on \(\mathbf{X}\) .

DEFINITION 2. A topological space is a set endowed with a topological structure.

The elements of a topological space are often called points. When a topology has been defined on a set X, this set is said to be the set underlying the topological space X.

Axiom \((\mathrm{O}_{\mathrm{I}})\) implies in particular that the union of the empty subset of D, i.e.~the empty set belongs to D. Axiom \((\mathrm{O}_{\mathrm{II}})\) implies that the intersection of the empty subset of D, i.e.~the set X, belongs to D.

Examples of topologies. Given any set X, the set of subsets of X consisting of X and \(\varnothing\) satisfies axioms \((\mathrm{O}_{\mathrm{I}})\) and \((\mathrm{O}_{\mathrm{II}})\) and therefore defines a topology on X. So does the set \(\mathfrak{P}(\mathbf{X})\) of all subsets of X: the topology it defines is the discrete topology on X, and the set X with this topology is called a discrete space.

A covering \((\mathbf{U}_i)_{i\in I}\) of a subset \(\mathbf{A}\) of a topological space \(\mathbf{X}\) is said to be open if all the \(\mathbf{U}_i\) are open in \(\mathbf{X}\) .

DEFINITION 3. A homeomorphism of a topological space X onto a topological space X' is an isomorphism of the topological structure of X onto that of X'; that is to say, in accordance with the general definitions a bijection of X onto X' which transforms the set of open subsets of X into the set of open subsets of X'.

X and \(X'\) are said to be homeomorphic if there is a homeomorphism of X onto \(X'\) .

Example. If X and \(X'\) are two discrete spaces, any bijection of X onto \(X'\) is a homeomorphism.

The following criterion follows immediately from the definition of a homeomorphism: for a bijection \(f\) of a topological space \(\mathbf{X}\) onto a topological space \(\mathbf{X}'\) to be a homeomorphism, it is necessary and sufficient that the image under \(f\) of each open set in \(\mathbf{X}\) is an open set in \(\mathbf{X}'\) , and that the inverse image under \(f\) of each open set in \(\mathbf{X}'\) is an open set in \(\mathbf{X}\) .

